\documentclass[12pt]{report}
\usepackage{graphicx} % Required for inserting images
\usepackage{subfigure}
\usepackage{cite}
\usepackage{geometry}
\usepackage{color}
\usepackage{amsmath}
\usepackage{amssymb}
\setcounter{secnumdepth}{4}
\setcounter{tocdepth}{4}
\usepackage{titlesec}
\linespread{1.6} 
% \setlength{\parindent}{0pt}
\usepackage{parskip}
\usepackage{multirow}
\usepackage[colorlinks=false, hidelinks]{hyperref}

\geometry{left=2cm,right=2cm,top=2cm,bottom=2cm}

\title{Post-doctoral: Simulation and measurement of road surface optical properties}
\author{Ruqin Xiao}
\date{This work is under \textbf{REFLECTIVITY} project}

  
\begin{document}

\maketitle
\begin{abstract}

    \textbf{Abstract--}
The spectral reflectance properties of road surfaces play a critical role in enhancing lighting efficiency, reducing environmental impacts, and supporting the development of advanced mobility systems. 
In this work, we present a novel method for measuring the spectral BRDF of road surfaces across a broad wavelength range (350–2500 nm). 
The approach is based on the gonioreflectometer developed by Cerema for photometric BRDF measurements, adapted for spectral analysis. 
Validation of the method is performed first using a Spectralon reference material with known diffuse reflectance properties, and then results are cross-validated with measurements conducted by Université Gustave Eiffel.
To illustrate the method's applicability, two road formulations (F2 and F3) and their individual components, provided by Spie Batignolles, were analyzed as representative examples. 
A total of 1728 measurements were performed for each sample, covering a wide range of illumination and observation geometries. 
The results show that the type of binder has a pronounced effect on the BRDF behavior of multi-component road surfaces, whereas aggregate size contributes only minor spectral variations, primarily in the specular reflection direction.
These findings could provide a comprehensive dataset and a methodological framework for future research on modeling the BRDF of multi-material road surfaces. 
Such models can combine the spectral characteristics of individual components through linear or nonlinear approaches. 



    \textbf{Key words--}
    Spectral BRDF, reflectance, solar albedo, reflectance of road surface, 
    BRDF models,
    BRDF measurements, 
    gonioreflectometer, goniospectroradiometer,
\end{abstract}

\tableofcontents

\input{introduction.tex}

\input{optical-properties-of-road-surface.tex}

\input{brdf-models-road.tex}

\input{measurement-literature.tex}

\input{measurement-road-sample.tex}

% \input{modelling-road-sample.tex}

\chapter{Conclusion and Perspectives}
%%%%%%%%%%%%%%%%%%%%%%%%%%%%%%%%%%%%%%%%%%%%%%%%%%%%%%%%%%%%%%%%%%%%%%%%%%%%%%%%
\section{Conclusion}
Both validations using spectralon with albedo of $0.5$ and comparing measured data to the ones measured by UGE, prove that the proposed spectral BRDF measurement approached shows a good accuracy.
For the study of spectral BRDF of roads surface, we conduct many measurements on the multi-material roads and their components.
The data demonstrates that the binder has a significant contribution on the BRDF behavior of the mixture roads.
Besides, the aggregates with different size has few difference except at the specular direction, which may have different contribution on the BRDF of mixture at the specular direction.


%%%%%%%%%%%%%%%%%%%%%%%%%%%%%%%%%%%%%%%%%%%%%%%%%%%%%%%%%%%%%%%%%%%%%%%%%%%%%%%%
\section{Perspectives}

In future work, the measured spectral BRDF data can be used to calibrate and evaluate existing reflectance models, such as the Cook Torrance model~\cite{1982_Cook}, the Oren Nayar model~\cite{1995_Oren}, and the Ashikhmin Shirley model~\cite{2000_Ashikhmin}. 
Furthermore, it will be of interest to develop a hybrid model that combines the most appropriate BRDF models for each material component to represent the BRDF of composite road surfaces.
This can provide a foundation for predicting the spectral BRDF of more complex pavement materials.

\bibliographystyle{unsrt} 
\bibliography{biblio}

\end{document}